% ============================================================================
% Chapter 11 --- Endpoint management
% ----------------------------------------------------------------------------
% Purpose : MDM, posture, BYOD, with explicit acknowledgement of OSS gap
%           on iOS/Android at scale.
%           Belongs to the Subject and Trust plane (D5.1).
% Page target: ~4 pages.
% Dependencies: Ch. 5 (DDM/DEP, MS-MDM, declarative Linux config),
%               Ch. 6 (identity attaches subjects to endpoints),
%               Ch. 7 (policy decisions consume posture claims),
%               Ch. 8 (network attachment depends on posture verification).
% Mirrors layer chapter template established in Ch. 6 (D6.5).
% ----------------------------------------------------------------------------
% Decisions registered in _coherence-EN.md:
%   D11.1 minimum compliance set for the endpoint layer
%   D11.2 reference instantiation = Munki/Ansible/osquery + MicroMDM (with gap)
%   D11.3 alternatives = Jamf+Intune; Workspace ONE
%   D11.4 OSS gap on iOS/Android at scale acknowledged transparently
%   D11.5 endpoint sovereignty scores
% ============================================================================

\chapter{Endpoint management}
\label{ch:endpoint}

\begin{caveatbox}[Dependencies]
\noindent This chapter applies the layer template of
\trchap{ch:identity}{} to the endpoint layer. It produces posture
claims consumed by \trchap{ch:policy}{}, contributes attributes
that \trchap{ch:identity}{} carries on session assertions, and is
constrained at the network attachment level by
\trchap{ch:network}{}. The layer belongs to the \emph{Subject and
trust plane}.
\end{caveatbox}

The endpoint layer covers institutional and personal devices that
connect to institutional services: managed Linux servers and
workstations, managed macOS, managed Windows, mobile devices
(\institution{iOS} and \institution{Android}), and BYOD endpoints
that may run any of the above. The layer reports \emph{posture}
claims to the policy layer, applies \emph{configuration} from a
declarative source, and exercises \emph{lifecycle controls}
(enrolment, audit, retirement) on the devices it manages.

\begin{caveatbox}[Honest acknowledgement]
\noindent The endpoint layer is the place where the open-source
reference instantiation is weakest, particularly for iOS and Android
at scale. This Reference does not minimise that gap; it documents
it (\S\ref{sec:ep-sovereignty}), proposes an open instantiation
where it is realistic, and treats the commercial alternatives as
the working choice for the segments where the open ecosystem has
not yet matured.
\end{caveatbox}

%-----------------------------------------------------------------------------
\section{Functional role and interface contract}
\label{sec:ep-contract}

\begin{contractbox}[Endpoint layer]
\noindent A compliant endpoint layer, in the sense of this
architecture, must:

\begin{itemize}[leftmargin=1.5em,nosep]
  \item enrol managed devices through a documented procedure that
    binds the device to an institutional identity;
  \item apply configuration from a declarative source: Apple
    declarative device management (DDM/DEP) on Apple platforms,
    MS-MDM on Windows, declarative configuration management
    (Ansible idempotent profiles, or equivalent) on Linux;
  \item produce posture claims consumed by \trchap{ch:policy}{}, in
    a schema that downstream policy authoring can rely on;
  \item support a documented decommissioning procedure (data wipe,
    enrolment removal, account deactivation) executable on demand;
  \item integrate with \trchap{ch:observability}{} so that
    endpoint events (enrolment, posture violation, configuration
    failure, suspicious activity) reach the institutional audit
    pipeline.
\end{itemize}

\noindent The contract does not require a single product to cover all
endpoint operating systems. Multi-product instantiations are
admissible provided that the institution operates a unified posture
schema across them, so that \trchap{ch:policy}{} can author rules
without per-OS branching for routine decisions.
\end{contractbox}

\paragraph{BYOD scope.} The contract applies to managed endpoints.
BYOD endpoints carry a more limited posture surface (no agent,
network-side observation only) and are governed through the
attachment policy of \trchap{ch:network}{} rather than through this
layer's enrolment procedures. The architectural principle is that
BYOD endpoints reach institutional services only through
contexts whose posture requirements they can satisfy without
agent-based verification.

%-----------------------------------------------------------------------------
\section{Reference instantiation: Munki + osquery + Ansible + MicroMDM}
\label{sec:ep-reference}

\begin{instantiationbox}[Endpoint layer]
\noindent\textbf{Topology by operating system.}
\textbf{Linux} endpoints (servers and workstations) are configured
through Ansible~\cite{ansible} idempotent playbooks under version
control; posture is queried through osquery~\cite{osquery}; package
state is tracked against an institutional repository.
\textbf{macOS} endpoints are configured through Munki~\cite{munki}
for software distribution and through MDM commands delivered by an
open-source MDM server (MicroMDM~\cite{micromdm} or NanoMDM, or a
comparable project) speaking Apple's MDM protocol with DEP/ADE
enrolment, and Declarative Device Management (DDM) where the chosen
open server and its companion services support it --- a coverage
currently partial relative to commercial MDMs; posture is
queried through osquery. \textbf{Windows} endpoints are partially
covered by declarative configuration through Ansible with the
WinRM/PowerShell-DSC integration, with the explicit observation
that the open-source coverage of MS-MDM is, at the time of writing,
substantially less mature than the Apple and Linux equivalents.
\textbf{iOS and Android} are addressed through commercial
alternatives in this architecture (\S\ref{sec:ep-alt1}); the
open-source ecosystem at the institutional scale would not, at the
time of writing, provide a credible reference instantiation.
\end{instantiationbox}

\paragraph{Posture schema.} The unified posture schema is the
artefact that decouples policy authoring from per-OS endpoint
products. The schema covers, at minimum: device identity,
enrolment status, OS version, encryption state, last-update
timestamp, presence of expected security agents, deviation from
declared configuration. Each product reports into this schema; the
mapping from product-specific signals to the unified schema is the
institution's responsibility.

\paragraph{Operational considerations.} A multi-product reference
instantiation imposes a multi-skill operating team. The indicative
staffing envelope set in the resourcing profile of \loc{institution-type}
spreads a small group
of engineers across the OS-specific operations, with one
engineer-equivalent dedicated to the unified posture schema and
its integration with the policy layer.

\begin{realworldbox}[Endpoint management precedents]
\noindent Research-intensive institutions are reported to operate
mixed-OS endpoint estates that combine open-source and commercial
tooling along the OS boundary, typically retaining open tooling for
Linux and macOS where the OSS ecosystem is mature, and commercial
tooling for Windows at scale and for iOS/Android. The specific tool
choices and operational details would be best confirmed against
the institution's own published documentation, when available, or
its presentations in the community fora of the institution's
national research and education network.
\end{realworldbox}

%-----------------------------------------------------------------------------
\section{Alternative~1 --- Jamf (macOS/iOS) + Intune (Windows) with OpenTelemetry export}
\label{sec:ep-alt1}

\begin{alternativebox}[Commercial multi-product]
\noindent The institution adopts Jamf~\cite{jamf} for Apple platforms
(macOS and iOS), Microsoft Intune~\cite{intune} for Windows and
Android, and continues Ansible/osquery for Linux. Posture data from each product is
exported through OpenTelemetry collectors so that the unified
posture schema is fed without any product becoming the
authoritative store.
\end{alternativebox}

\paragraph{Where it adds value.} For iOS and Android in particular,
the operational depth of Jamf and Intune is, at the time of
writing, substantially ahead of the open alternatives. For an
institution where the mobile estate is large and lifecycle controls
must be exercised reliably, this is the working choice.

\paragraph{Where it constrains the institution.} Each product
expresses configuration and posture in its own model; the
reconciliation of these models with the unified schema of
\S\ref{sec:ep-reference} is a sustained engineering effort. Beyond
the schema work, the institution accepts a non-trivial financial
commitment per device per year, and a roadmap that the suppliers
control. The architectural posture is to keep posture data
exported through OpenTelemetry collectors so that the policy layer
remains the source of authorisation truth, and so that exit from
either product would be a re-implementation of the data pipe rather
than a re-foundation of the policy.

\paragraph{When it could be the right choice.} An institution
whose iOS and Android estates are operationally significant, whose
financial position absorbs the per-device licensing, and whose
sovereignty posture accepts a Partial score on the technical and
institutional dimensions of the endpoint layer in exchange for the
operational depth.

%-----------------------------------------------------------------------------
\section{Alternative~2 --- Workspace ONE multi-OS}
\label{sec:ep-alt2}

\begin{alternativebox}[Single-vendor multi-OS]
\noindent The institution adopts a single multi-OS endpoint
management platform (Workspace ONE, or comparable) covering Apple,
Windows, Android and a subset of Linux through one supplier
relationship. Posture export to OpenTelemetry remains a contractual
requirement.
\end{alternativebox}

\paragraph{Where it adds value.} A single supplier relationship,
a single skill set within the team, and consistent posture semantics
across operating systems out of the box. The reconciliation work
that the multi-product alternative imposes is largely absorbed by
the supplier.

\paragraph{Where it constrains the institution.} The single supplier
becomes the sole point of failure for the entire endpoint layer.
Substitutability is significantly weaker than in the multi-product
alternative, and institutional negotiation leverage diminishes
across renewal cycles. The exit-cost-proportionality indicator of
\trchap{ch:sovereignty-grid}{} should be evaluated explicitly at
procurement.

\paragraph{When it could be the right choice.} An institution with
limited engineering capacity to integrate multiple endpoint
products, where the operational simplification is more valuable
than the substitutability cost, and whose financial sovereignty
posture is comfortable with a single-supplier exposure on the
endpoint layer.

%-----------------------------------------------------------------------------
\section{Sovereignty grid applied}
\label{sec:ep-sovereignty}

\begin{table}[H]
\centering\small
\caption{Per-dimension sovereignty profile of the endpoint
instantiations. The reference is scored across two columns to
make the iOS/Android gap visible.}
\label{tab:ep-sovereignty}
\begin{tabularx}{\textwidth}{%
  >{\hsize=0.6\hsize\linewidth=\hsize\bfseries\raggedright\arraybackslash}X%
  >{\hsize=1.2\hsize\linewidth=\hsize\centering\arraybackslash}X%
  >{\hsize=1.2\hsize\linewidth=\hsize\centering\arraybackslash}X%
  >{\hsize=1.0\hsize\linewidth=\hsize\centering\arraybackslash}X%
  >{\hsize=1.0\hsize\linewidth=\hsize\centering\arraybackslash}X%
}
\toprule
Dimension & Reference \emph{Linux/macOS} & Reference \emph{iOS/Android} & Jamf/Intune & Workspace ONE \\
\midrule
Data          & 3 Robust      & --- (gap)     & 2 Established & 2 Established \\
Technical     & 3 Robust      & 0--1 Absent/Partial & 1 Partial     & 1 Partial     \\
Financial     & 2 Established & 0--1 Absent/Partial & 1 Partial     & 1 Partial     \\
Operational   & 2 Established & 0--1 Absent/Partial & 3 Robust      & 3 Robust      \\
Institutional & 3 Robust      & 0--1 Absent/Partial & 1 Partial     & 1 Partial     \\
\bottomrule
\end{tabularx}
\end{table}

\noindent The transparent split of the reference column is
deliberate. The open-source reference is competitive on Linux and
macOS; it is not, at the time of writing, a credible institutional
choice for iOS and Android. The architectural posture is to use the
reference where it works and the commercial alternatives where the
ecosystem requires it, while keeping the policy layer in a position
to consume posture from any of them without favour.

\paragraph{Driving indicators by dimension.} Data: enrolment data,
posture data, telemetry destination. Technical: source visibility,
captive extensions in the configuration model. Financial: per-device
licensing and renewal trajectory. Operational: skills market depth
across the institution's OS mix. Institutional: roadmap influence
and integration freedom.

%-----------------------------------------------------------------------------
\section{Regulatory anchoring}
\label{sec:ep-regulatory}

\noindent The endpoint layer is the device-side anchor of the
applicable data protection regime's technical-measures obligation
(technical measures applied at the device that processes personal
data) and the principal locus of the applicable cybersecurity
baseline regime's requirements on basic cyber hygiene practices,
asset management, and indirectly authentication (where the endpoint
carries the authentication factors). Posture attestations consumable by the
policy layer operationalise the device dimension of
zero-trust~(SP~800-207). The layer contributes to 27001~A.8.1
(user endpoint devices), A.8.7 (protection against malware) and
A.8.9 (configuration management) and to NIST CSF \textsf{PROTECT}
(PR.PS --- platform security). Detailed mapping in
\trchap{ch:regulatory-mapping}{} \S\ref{sec:rm-gdpr},
\S\ref{sec:rm-nis2}, \S\ref{sec:rm-iso}, \S\ref{sec:rm-nist-csf}.

%-----------------------------------------------------------------------------
\section{Common pitfalls}
\label{sec:ep-pitfalls}

\begin{pitfallbox}[BYOD posture without enforcement]
\noindent The institution declares a BYOD posture (encryption,
patch level, screen lock) and relies on user attestation for it.
Within months, the actual fleet has drifted from the declared
posture, and the network attachment of \trchap{ch:network}{} no
longer reflects what the institution thought it enforced. The
mitigation is to make BYOD posture either verifiable
(network-side observation, captive-portal challenge) or
contractual without the architecture relying on it for sensitive
zones.
\end{pitfallbox}

\begin{pitfallbox}[MDM cost explosion]
\noindent Per-device licensing in the commercial alternatives can
grow with the estate at a rate that procurement did not anticipate,
particularly when contractor and visiting-researcher devices accrue.
The mitigation is to budget per-device licensing as a recurring
variable cost, to maintain a current device inventory, and to
include the projection in the financial-sovereignty review.
\end{pitfallbox}

\begin{pitfallbox}[Policy drift between MDM and central PDP]
\noindent An MDM product expresses authorisation rules in its own
language (compliance policies, conditional access expressions); the
central PDP of \trchap{ch:policy}{} expresses them in Rego or Cedar.
Over time, the two drift, and the same notional rule says different
things in the two systems. The mitigation is to author posture-aware
authorisation in the central PDP only, and to use the MDM
product's authorisation features only for device-level decisions
that the policy layer does not need to know about.
\end{pitfallbox}

\begin{pitfallbox}[Undocumented MDM extensions]
\noindent Commercial MDM products extend the standard
DDM/MS-MDM/declarative protocols with proprietary capabilities.
Adoption of such extensions in critical paths erodes
replaceability without making it visible. The mitigation is to
audit the configuration profiles and policies for non-standard
extensions periodically, and to confine such extensions to
non-critical paths or to substitute them.
\end{pitfallbox}
